\chapter*{Annexe}\addcontentsline{toc}{chapter}{Annexe}\label{annexe1}

Cette annexe est un \textbf{exemple}. Remplacez ce texte et ces
figures par vos propres contenus, i.e. les captures d'écran, schémas, tableaux de données,
extraits de configuration, ... . Les images utilisées ici sont de simples images de
substitution (\emph{placeholders}) à remplacer par les vôtres dans le dossier \texttt{images/}.

\subsection*{Insérer une figure}

Pour insérer une figure, placez votre image dans \texttt{images/} et référencez-la ensuite avec
\verb|\includegraphics|. Vous pouvez y renvoyer dans le texte, comme à la
figure~\ref{fig:exemple1}.

\begin{figure}[H]
	\centering
	\includegraphics[width=0.7\textwidth]{images/exemple-figure-1.png}
	\caption{Exemple de figure (image de substitution).}
	\label{fig:exemple1}
\end{figure}

\subsection*{Ajuster la taille}

La largeur se règle avec l'option \verb|width| (par exemple \verb|width=0.5\textwidth|)
ou avec \verb|scale|. La figure~\ref{fig:exemple2} occupe la moitié de la largeur du texte.

\begin{figure}[H]
	\centering
	\includegraphics[width=0.5\textwidth]{images/exemple-figure-2.png}
	\caption{Deuxième exemple de figure, en demi-largeur.}
	\label{fig:exemple2}
\end{figure}

\subsection*{Capture d'écran}

Les captures d'écran s'insèrent de la même manière. Voir la figure~\ref{fig:capture}.

\begin{figure}[H]
	\centering
	\includegraphics[width=\textwidth]{images/exemple-capture.png}
	\caption{Exemple de capture d'écran (image de substitution).}
	\label{fig:capture}
\end{figure}
